\section{Related Work}
\label{sec:related}

CSM-SAM sits at the intersection of four lines of work: promptable
foundation segmenters and their memory-augmented video extensions,
longitudinal medical segmentation with registration priors,
post-treatment tumor segmentation on BraTS, and bitemporal change
detection in remote sensing. We review each and then position the
contribution.

\subsection{Segment Anything family and memory-augmented variants}
\label{sec:related:sam}

SAM~\cite{kirillov2023sam} introduced promptable segmentation at
foundation scale: a ViT-H encoder, a prompt encoder, and a
lightweight mask decoder trained on SA-1B ($\sim$1B masks, 11M
images) via a model-in-the-loop data engine. It produces zero-shot
masks from a point, box, or coarse-mask prompt but has no temporal
context: each call operates on a single image in isolation.

SAM~2~\cite{ravi2024sam2} extends the design to video by adding a
memory bank and a memory-attention block. A streaming encoder writes
object-conditioned ``memory tokens'' for past frames into a bounded
queue; a memory-attention module inserted between the image encoder
and the mask decoder lets the current frame's tokens attend into that
queue through standard cross-attention, recovering object identity
across occlusion or motion gaps without re-prompting. The memory bank
is \emph{within-video}: all tokens come from the same continuous
acquisition.

MedSAM2~\cite{ma2024medsam2} adapts SAM~2 to 3D medical volumes by
treating the slice stack as a short video. A prompt on one slice
propagates through the volume via the same memory-attention block,
yielding 3D masks from 2D prompts with decoder-only fine-tuning. The
memory bank is still within-session clinically: all slices belong to
a single acquisition of one patient at one visit.

CSM-SAM continues this trajectory by introducing a \emph{cross-session}
memory. The tokens written into the bank come from a \emph{different
scan of the same subject} separated by one to four weeks of treatment,
rather than neighbouring slices of one volume. The memory-attention
block that SAM~2 uses to bridge temporal gaps within a video is reused
verbatim to bridge the much larger inter-session gap between pre-RT
(week 0) and mid-RT (week 2-3) scans, with two additions: an
XOR-derived change-map head and a continuous weeks-elapsed embedding.
SAM, SAM~2, and MedSAM2 do not expose this cross-session mode.

\subsection{Longitudinal medical image segmentation}
\label{sec:related:longi}

In single-timepoint medical segmentation, 3D specialists dominate.
nnU-Net~\cite{isensee2021nnunet} is a self-configuring pipeline
around a volumetric U-Net that chooses preprocessing, patch size, and
training schedule from a dataset fingerprint; trained from scratch on
a few hundred volumes with heavy augmentation and five-fold
ensembling, it remains the reference baseline on most medical
benchmarks. Swin\,UNETR~\cite{hatamizadeh2021swinunetr} replaces the
CNN encoder with a 3D Swin Transformer;
UNETR~\cite{hatamizadeh2022unetr} uses a plain ViT encoder in the
same U-Net layout; MedNeXt~\cite{roy2023mednext} is a ConvNeXt-style
3D U-Net with depthwise convolutions and kernel up-scaling init. All
four train end-to-end at 30-90M parameters with no built-in mechanism
for consuming a prior visit.

The HNTS-MRG 2024 challenge~\cite{wahid2024hntsmrg} formalizes the
adaptive-radiotherapy setting CSM-SAM targets: given pre-RT (week 0)
T2-weighted MRI with ground-truth masks, segment the mid-RT
(week 2-3) scan of the same patient. The organizers ship a
\texttt{preRT\_mask\_registered} field -- the pre-RT mask warped onto
the mid-RT grid by deformable B-spline registration -- precisely
because the task is longitudinal. Every published top-5 entry
consumes this registered mask as an auxiliary input channel. The top
team (UW\,LAIR, aggDSC~0.733) uses a
SegResNet~\cite{myronenko2019segresnet} backbone with a mask-aware
attention module that modulates decoder features by the warped pre-RT
mask; teams 2-5 (0.727-0.710) are nnU-Net~3d\_fullres variants with
the warped mask stacked on the input.

The structural property these pipelines share is that pre/mid
correspondence is \emph{engineered upstream}. Classical deformable
registration produces the warped mask once, offline; the segmenter
treats it as a static prior channel and never revisits the raw pre-RT
image. This decouples correspondence from segmentation and exposes
the system to registration error in low-contrast soft tissue --
precisely the regime (post-chemoradiation head-and-neck tumor) where
B-spline alignment is least reliable. It also discards the pre-RT
intensity information: only the binary mask, not the pre-RT image
itself, reaches the decoder. CSM-SAM instead learns pre/mid
correspondence end-to-end through cross-time token attention on the
frozen SAM~2 encoder grid, conditions this attention on weeks elapsed,
and supervises it with the XOR-derived change map. No deformable
registration is invoked at training or inference time.

\subsection{BraTS and post-treatment tumor segmentation}
\label{sec:related:brats}

The BraTS 2024 post-treatment glioma track~\cite{delarosa2024brats}
is the other major longitudinal medical benchmark: paired baseline
(TP000) and follow-up (TP001) multi-parametric MRI with four
co-registered sequences (T1c, T1n, T2w, T2f/FLAIR), evaluated by
lesion-wise Dice on four sub-regions. Top solutions are nnU-Net,
MedNeXt, and Swin\,UNETR 3D ensembles at 128$^3$ patches, each
trained as a four-channel-in, multi-class-out volumetric segmenter.
The critical observation -- explicit in the overview paper -- is that
these pipelines process TP000 and TP001 \emph{independently}: a
single model runs on each visit with no mechanism for conditioning
the follow-up prediction on the baseline image tokens or mask.
Longitudinal structure is unused at model level and resurfaces only
in clinical post-processing. CSM-SAM's cross-session memory occupies
this gap: TP000 tokens are written into a memory bank and TP001
tokens attend into it, directly conditioning the follow-up prediction
on the baseline encoding.

\subsection{Change detection in remote sensing}
\label{sec:related:cd}

Remote-sensing change detection has developed a parallel bitemporal
toolbox whose architectural choices mirror -- but do not anticipate
-- CSM-SAM's. FC-Siam~\cite{daudt2018fcsiam} introduced the
fully-convolutional siamese U-Net: two weight-shared encoder branches
process pre- and post-event tiles, and the decoder fuses
skip-connection features by concatenation (FC-Siam-conc) or absolute
difference (FC-Siam-diff) at every scale, establishing the template
subsequent binary CD architectures follow. SNUNet-CD~\cite{fang2022snunet}
replaces the U-Net with a NestedUNet and adds an Ensemble Channel
Attention Module (ECAM) over nested-skip features.

Transformer-based CD moved fusion out of the CNN skips and into
token space. BIT~\cite{chen2021bit} uses a ResNet-18 siamese CNN to
extract bitemporal feature maps, projects each into a small set of
``semantic tokens'' (4-8 per image), and refines them with a shared
transformer encoder-decoder before reprojecting to pixel space.
ChangeFormer~\cite{bandara2022changeformer} goes all-transformer: a
hierarchical MiT encoder operates siamesely on both tiles and an MLP
decoder fuses per-scale feature differences.
TinyCD~\cite{codegoni2023tinycd} compresses the recipe into
$\sim$0.3M parameters with a truncated EfficientNet siamese encoder
and a Mix-and-Attention Mask Block.
ChangeMamba~\cite{chen2024changemamba} replaces self-attention with
Visual State-Space (Mamba) blocks and scans the stacked bitemporal
feature tensor along multiple raster directions for linear-time
long-range mixing.

The line closest to CSM-SAM in spirit is SAM-CD~\cite{ding2024samcd},
which freezes a SAM~v1 ViT encoder, attaches a FastSAM-style adapter
and a change decoder, and trains with a semantic-consistency loss
encouraging feature similarity between spatially overlapping
non-changed regions across timestamps. It demonstrates that a frozen
SAM encoder transfers to CD despite never seeing a bitemporal signal
during pretraining. However, SAM-CD uses the original SAM (no
memory-attention machinery), couples timestamps via a per-image
feature-similarity loss rather than cross-time token attention, and
has neither a change-map auxiliary head nor a temporal-gap embedding.
CSM-SAM uses the more capable SAM~2 backbone, reuses its
memory-attention block at the cross-session level, and adds both
signals.

The common thread across FC-Siam, SNUNet-CD, BIT, ChangeFormer,
TinyCD, and ChangeMamba is a from-scratch-or-ImageNet backbone
trained on small CD corpora, with bitemporal fusion implemented by
differencing, concatenation, nested skips, low-rank transformer
tokens, or directional SSM scans -- not by a memory-bank read.
CSM-SAM replaces the small from-scratch backbone with a frozen
1B-mask SAM~2 ViT and replaces the hand-designed fusion operator
with the memory-attention mechanism SAM~2 uses for within-video
persistence, adapted to the cross-session regime.

\subsection{Semantic change detection}
\label{sec:related:scd}

Semantic change detection is the structurally closest analog to
CSM-SAM: the task produces the full $\{\text{pre-mask},
\text{post-mask}, \text{change-map}\}$ triple, exactly matching
CSM-SAM's output signature. HRSCD~\cite{daudt2019hrscd} introduces a
dual-branch multi-task FCN in which two semantic branches share
encoder features with a dedicated change branch; its ``strategy~4''
variant trains all three heads jointly with a consistency loss
penalizing semantic pairs incompatible with the predicted change.
Bi-SRNet~\cite{ding2022bisrnet} pairs a siamese ResNet encoder with
a Siamese Relation-aware (SR) module for intra-temporal spatial
context and a Cross-Temporal Relation module for inter-temporal
exchange, feeding three heads (pre-semantic, post-semantic, binary
change) with a semantic-consistency loss.
SCanNet~\cite{ding2024scannet} extends this with a Semantic Change
Transformer that attends symmetrically over concatenated bitemporal
tokens, modelling ``pixel at date~1 vs pixel at date~2'' dependencies
via three decoding heads and a joint semantic-change loss; it is
currently the most widely reproduced SECOND result.

Two properties of this family matter for positioning. First, the
architectures are specialized siamese CNNs or CNN+transformer hybrids
of 20-40M parameters trained from scratch on SECOND's 6-class urban
land-cover label set; they do not transfer to medical tasks or to
disaster-damage grading without architectural re-engineering and
retraining. Second, bitemporal coupling is symmetric -- both dates
are treated equivalently and feature exchange flows both ways.
CSM-SAM's cross-session memory is asymmetric by design: the pre-RT
visit writes a read-only memory bank that the mid-RT query attends
into, mirroring the clinical workflow in which the pre-RT scan is
always available first. The same module, unchanged, produces the
$\{\text{pre}, \text{post}, \text{change}\}$ triple on top of a frozen
foundation backbone never trained on land-cover or medical data.

\subsection{Our positioning}
\label{sec:related:positioning}

Relative to the four lines above, CSM-SAM occupies a corner none of
them covers jointly. It (i) builds on a frozen foundation backbone
(SAM~2 ViT-H, pretrained on SA-1B masks and SA-V video, no medical
fine-tuning of the encoder) instead of a from-scratch medical or CD
specialist; (ii) reuses SAM~2's within-video memory-attention block
as a \emph{cross-session} mechanism, learning pre/mid correspondence
end-to-end rather than outsourcing it to deformable registration
(HNTS-MRG winners) or symmetric siamese fusion (CD literature);
(iii) adds an auxiliary change-map head supervised by the pre/mid
mask XOR, a free signal every HNTS-MRG baseline throws away and no
SAM/SAM~2 variant exposes; and (iv) conditions the memory attention
on a continuous weeks-elapsed embedding, calibrating expected change
magnitude against the 1-4-week treatment-response window. The
trainable parameter count is $\sim$2M (memory module plus final
decoder layer), and the same module transfers unchanged from
HNTS-MRG to post-treatment BraTS, to SECOND, and to xBD, which is
the generality claim the subsequent sections evaluate.
